\subsection{Avoiding retrieval when the decision is already certified}
\label{sec:retrieval-gate}
The same affine loss bounds can provide a stopping rule for a downstream
model-choice request. Here the request is to return a model whose block MSE is
within $\epsilon=0.01M_2$ of the best candidate, where $M_2$ is the previously
fixed training-period second moment. This differs from the joint half-width
requirement. With retained block energy $E$, we stop before retrieval whenever
\[
 \min_m\max\left\{0,\max_{j\ne m}
 \left[d_{mj}+E\max_h c_{mj,h}\right]\right\}\leq\epsilon.
\]
Every permitted completion then satisfies the requested tolerance for a
qualifying model. Otherwise, one bulk fetch supplies the five internal counter
values used by all four models. We use the first qualifying model in
lexicographic order and cache every returned value.

On the existing panels, this rule avoids 8 of 819 original and 198 of 279
external block fetches when endpoints are already retained: request reductions
of 0.98\% and 70.97\%, respectively. The per-system counts are $0,1,7$ of
273 original blocks and $75,37,86$ of 93 external blocks. Exact arithmetic
checks all 39,528 loss differences at the simplex vertices. A fixed-order
timing pass takes 0.265\,s for the 1,098 gates after input preparation
(median 0.265\,ms per block); loading and rational conversion take a further
0.490\,s on the reported machine.

\begin{table}[htbp]
\centering\small
\caption{Range-access replay without later forecast copies:
requests/returned counter values. One fetch is
shared across all models. Cold access charges endpoint acquisition; period
access coalesces ranges per system.}
\label{tab:retrieval-gate-cost}
\begin{tabular}{llrr}
\toprule
Access & Policy & Original & External\\
\midrule
Immediate, endpoints held & Always bulk & 819/4095 & 279/1395\\
                         & Certificate gate & 811/4055 & 81/405\\
Immediate, cold          & Direct bulk & 819/5187 & 279/1767\\
                         & Endpoints, then gate & 1630/5147 & 360/777\\
Whole period, cold       & Direct bulk & 3/6537 & 3/2217\\
                         & Coarse grid, then gated ranges & 6/7623 & 6/2403\\
\bottomrule
\end{tabular}
\end{table}

For a cold audit priced per request, the router selects direct bulk. In a
historical audit it first uses permitted forecast copies, with binary64
rounding enclosures: these certify 810 original and 270 external blocks.
With endpoints retained, after the initial gate only nine original and three external blocks need
fine retrieval, covered by three requests/51 values and two requests/22 values.

Table~\ref{tab:retrieval-gate-cost} replays the declared access schedules.
We separately validated the interface with 12 predeclared public API requests,
including full-record downloads. They confirm the range semantics, all 279
external totals, and the tolerance of all 198 decisions selected before
retrieval. The measured period-wide coarse-plus-gated requests receive 69,809
response-body bytes in six requests, versus 64,357 in three bulk requests.
The 70.97\% request reduction therefore belongs to immediate audits with
retained endpoints. The comparison measures requests and returned data under
one fixed decision tolerance; it supplies a concrete stopping rule rather
than a price or latency estimate.

\begin{table}[htbp]\centering\small
\caption{Absolute tolerance and realized loss scales on the same blocks. All loss columns are in kWh$^2$. Med. MSE pools the four fixed models across blocks; best and spread are blockwise minimum and maximum-minus-minimum MSE, respectively, before taking medians. All four counts blocks whose realized spread is at most $\epsilon$; it is not a certificate count.}
\label{tab:retrieval-tolerance-scale}
\begin{tabular}{lrrrrr}\toprule
System & $\epsilon$ & Med. MSE & Med. best & Med. spread & All four\\\midrule
BLEM 1 & 0.003575 & 0.02371 & 0.01122 & 0.03286 & 12/273\\
BLEM 2 & 0.004427 & 0.03981 & 0.01881 & 0.05139 & 15/273\\
BLEM 3 & 0.003251 & 0.05632 & 0.02384 & 0.08984 & 0/273\\
HPM 1 & 0.006101 & 0.02327 & $4.292\times10^{-6}$ & 0.1641 & 0/93\\
HPM 2 & 0.002245 & 0.002964 & 0.000354 & 0.005857 & 25/93\\
HPM 7 & 0.009255 & 0.0004555 & $7.2\times10^{-5}$ & 0.009809 & 44/93\\
\bottomrule\end{tabular}\end{table}
The tolerance has different implications across systems. HPM~7 has all four models within tolerance on 44 of 93 blocks, while HPM~1 has none and has a near-zero median best-model error alongside a much larger model spread. The blockwise best MSE is a hindsight reference, not the performance of one deployable model. These summaries interpret the unchanged tolerance rather than calibrating it to an operational loss.
